\section{Responsibility and duty: you have the power to prevent this}
\label{sec:responsibility}

The fifth appeal is addressed to the legislator directly. It is not sympathy but
obligation: you have the power to prevent this, and future generations will judge
what we do now. This is the moment the narrative turns from describing a problem
to naming the body that can solve it and the precise instrument, H. R. 9510, that
lets it.

[DRAFTING INSTRUCTIONS: in the full-narrative, frame the duty appeal with
\texttt{review/references/narrative\_refs.bib} entry \texttt{moreland2015hearing}
(how credible testimony informs the legislative record) and weave in the MAIN
ACTIONS (advocacy, coalition building, policy entrepreneurship, reconciliation)
and the KEY TERMS (testimony and committee hearings, markup). Reproduce Mermaid
figures \texttt{review/mermaid/diagrams/03-bill-to-law-process.md} and
\texttt{14-reconciliation-gitgraph.md} as colored TikZ. Add a body-width table
mapping each MAIN ACTION to the committee step where it has effect.]

\begin{narrativefig}
\node[nfone] (a) at (0,0) {Introduce};
\node[nfact] (b) at (3.4,0) {Markup};
\node[nftwo] (c) at (7.0,0) {Reconcile};
\node[nfhope] (d) at (10.2,0) {Enact};
\draw[nfedge] (a)--(b); \draw[nfedge] (b)--(c); \draw[nfedge] (c)--(d);
\end{narrativefig}
\vspace{-0.2cm}
\figcaption{Figure 6. The duty to move a bill from introduction to enactment
(placeholder; expanded from Mermaid 03 in the full-narrative).}

Responsibility is the bridge from feeling to action. It tells the committee that
the next move is theirs and that the record will remember it.
